\documentclass[11pt]{article}


\usepackage{fullpage}
\usepackage{graphicx}
\usepackage{amsmath}
\usepackage{amssymb}
\usepackage{amsthm}
\usepackage{fancyvrb}

\newcommand{\myname}{Mehshan Mustafa}

\newenvironment{theorem}[2][Theorem]{\begin{trivlist}
\item[\hskip \labelsep {\bfseries #1}\hskip \labelsep {\bfseries #2.}]}{\end{trivlist}}
\newenvironment{lemma}[2][Lemma]{\begin{trivlist}
\item[\hskip \labelsep {\bfseries #1}\hskip \labelsep {\bfseries #2.}]}{\end{trivlist}}
\newenvironment{exercise}[2][Exercise]{\begin{trivlist}
\item[\hskip \labelsep {\bfseries #1}\hskip \labelsep {\bfseries #2.}]}{\end{trivlist}}
\newenvironment{problem}[2][Problem]{\begin{trivlist}
\item[\hskip \labelsep {\bfseries #1}\hskip \labelsep {\bfseries #2.}]}{\end{trivlist}}
\newenvironment{question}[2][Question]{\begin{trivlist}
\item[\hskip \labelsep {\bfseries #1}\hskip \labelsep {\bfseries #2.}]}{\end{trivlist}}
\newenvironment{corollary}[2][Corollary]{\begin{trivlist}
\item[\hskip \labelsep {\bfseries #1}\hskip \labelsep {\bfseries #2.}]}{\end{trivlist}}
\newenvironment{solution}{\begin{proof}[Solution]}{\end{proof}}
\newenvironment{idea}[2][Proof Idea.]{\textit{#1} #2}



\parindent0in
\pagestyle{plain}
\thispagestyle{plain}

\usepackage{csquotes}
\usepackage[shortlabels]{enumitem}

\newcommand{\dated}{\today}
\newcommand{\token}[1]{\langle \text{#1} \rangle}

\begin{document}

\textbf{Introduction to the Theory of
Computation}\hfill\textbf{\myname}\\[0.01in]
\textbf{Chapter 8: Space Complexity}\hfill\textbf{\dated}\\
\smallskip\hrule\bigskip

\begin{problem}{8.30}
Show that $E_{\text{DFA}}$ is NL-complete.
\end{problem}

\begin{proof}
To show that $E_{\text{DFA}}$ is NL-complete, we must show that it is in NL and that all NL problems are log space reducible to it. To show $E_{\text{DFA}} \in$ NL, we show that $\overline{E_{\text{DFA}}} \in$ NL. We give an algorithm $M$ that decides $\overline{E_{\text{DFA}}}$ in non-deterministic log space. \\

$M =$ \textquotedblleft On input $\langle D \rangle$, where $D = (Q, \Sigma, \delta, q_0, F)$ is a DFA and every state $q_i \in Q$ is a binary integer greater than 0:
\begin{enumerate}
\item Let $i$, $k$, $n$ and $q$ be binary integers on work tape.
\item Let $n = |Q|$.
\item Non-deterministically, repeat for each $k = 0, 1, \dots, n$:
\item \hspace*{0.5cm} Let $q = q_0$ and $i = 0$.
\item \hspace*{0.5cm} If $i < k$, then:
\item \hspace*{1.0cm} Non-deterministically, repeat for each $a \in \Sigma$:
\item \hspace*{1.5cm} Let $q = \delta(q, a)$ and $i = i + 1$.
\item \hspace*{1.5cm} Go to step 5. 
\item \hspace*{0.5cm} If $q \in F$, then \textit{accept}. Otherwise, \textit{reject}.\textquotedblright
\end{enumerate}
\end{proof}
\end{document}